% TOP_PROBLEM: {"id": 1098, "title": "Bounds for Homogeneous Polynomial Zeros", "category_id": 1, "category": {"id": 1, "name": "number_theory", "display_name": "Number Theory", "description": "Properties of integers, prime numbers, Diophantine equations.", "slug": "number-theory", "order_index": 1, "created_at": "2026-07-31T15:26:25.670Z"}, "status": "open"}
% FIRST_POSTED: null
% ENTRY_KIND: statement_only
% Imported from: batch03.tex; UnsolvedMath ID 1098
\documentclass[11pt]{article}
\usepackage[margin=25mm]{geometry}
\usepackage{fontspec}
\setmainfont{Latin Modern Roman}
\usepackage{amsmath,amssymb,mathtools}
\usepackage{unicode-math}
\setmathfont{Latin Modern Math}
\usepackage{xurl,hyperref}
\hypersetup{colorlinks=true,urlcolor=blue}
\setcounter{secnumdepth}{0}
\setlength{\parindent}{0pt}
\setlength{\parskip}{5pt}
\setlength{\emergencystretch}{3em}
\newcommand{\sref}[2]{\hyperlink{source:#1:#2}{[S#2]}}
\newcommand{\cataloguescope}[1]{\paragraph{Recorded target.}#1}
\begin{document}
% UnsolvedMath ID 1098; source catalogue code GREEN-098
\setcounter{section}{1097}
\section{Bounds for Homogeneous Polynomial Zeros}\label{1098}
{\small\sffamily UnsolvedMath ID 1098 · Number Theory}
\subsection{Definitions and mathematical statement}

\paragraph{Shared notation.}$\mathbb N=\{1,2,\ldots\}$, and $\mathbb Z,\mathbb Q,\mathbb R,\mathbb C$ denote the integers, rationals, reals and complex numbers. Any other conventions are specified locally.


A polynomial $F\in\mathbb Z[x_1,\ldots,x_n]$ is a homogeneous form of degree $d$ if every monomial with nonzero coefficient has total degree $d$. A nontrivial integral zero is a vector $\mathbf x\in\mathbb Z^n\setminus\{\mathbf0\}$ with $F(\mathbf x)=0$. Such a zero exists exactly when a nonzero rational zero exists, by multiplying rational coordinates by a common denominator. For each odd $d\geq3$, define $\nu(d)$ as the least nonnegative integer $t$ such that every such form in every number $n>t$ of variables has a nontrivial integral zero; Birch's theorem ensures that $\nu(d)$ is finite.
\paragraph{Statement recorded in the supplied source.}
For every odd integer $d\geq3$, obtain explicit upper bounds for $\nu(d)$ and improve their dependence on $d$; that is, find quantitatively sharp functions $B(d)$ for which
\[
 n>B(d)\ \Longrightarrow\
 \bigl(\forall\text{ homogeneous }F\in\mathbb Z[x_1,\ldots,x_n]
 \text{ of degree }d\bigr)
 \bigl(\exists\mathbf x\in\mathbb Z^n\setminus\{\mathbf0\}\bigr)
 F(\mathbf x)=0.
\]
No nonsingularity, diagonal form, or coefficient-size restriction is imposed. Bounds are requested uniformly over all forms of the given odd degree. \sref{1098}{2}
\subsection{Short English statement}
Bound how many variables suffice to force a nonzero integer zero of every homogeneous polynomial of a given odd degree.
\subsection{Sources}
\begin{description}
\item[\hypertarget{source:1098:1}{S1}] ULAM.AI and the UnsolvedMath Contributors, \emph{UnsolvedMath: A Curated Collection of Open Mathematics Problems}, record 1098 (GREEN-098). CC BY 4.0. \url{https://huggingface.co/datasets/ulamai/UnsolvedMath}.
\item[\hypertarget{source:1098:2}{S2}] Amichai Lampert, Andrew Snowden and Tamar Ziegler, \emph{Polynomial Bounds for Birch’s Theorem}, arXiv:2512.00697v1 (2025), §1, definition of $N(d,s)$ and discussion of general degree bounds. For $K=\mathbb Q$, the recorded single-form threshold is $\nu(d)=N(d,1)-1$; the paper distinguishes fixed-degree dependence on the number of forms from dependence on degree. \url{https://arxiv.org/abs/2512.00697}.
\item[\hypertarget{source:1098:3}{S3}] Ben Green, \emph{100 Open Problems}, maintained author list, Problem 98 and its comments, accessed 7 October 2026. \url{https://people.maths.ox.ac.uk/greenbj/papers/open-problems.pdf}.
\end{description}
\label{end:1098}
\end{document}
